\documentclass{article}
\usepackage{anyfontsize}
\usepackage[UTF8]{ctex} % 如果需要支持中文
\usepackage{fontspec} 
\setCJKmainfont{SourceHanSansCN-Light}[
    BoldFont=SourceHanSansCN-Medium
]

\usepackage[a4paper, left=0.1in, right=0.1in, top=0.05in, bottom=0.05in]{geometry} % 设置四边页边距为0.1英寸{geometry} % 设置页边距为0.5英寸
\usepackage{enumitem} % 用于自定义列表
\usepackage{amsmath}  % 数学公式支持
\usepackage{tikz}
\usetikzlibrary{arrows.meta, positioning}
\setlength{\parindent}{0pt} % 不缩进
\usepackage{etoolbox} % 用于列表操作
%调整类代码
\usepackage{comment}
\usepackage{tocloft} % 用于定制目录 样式
\usepackage{hyperref} %不需要目录盒子注释这
\usepackage{ifthen}
\usepackage{tikz}
\usetikzlibrary{arrows.meta}
\usepackage{titletoc}
\usepackage{graphicx}
\usepackage{float}

\usepackage{amssymb}  % 提供数学符号，包括\therefore
%目录设定
%快捷键 CMD + / 注释
%  \renewcommand{\cftdot}{} % 隐藏目录中的页码
%  \cftpagenumbersoff{section}
%  \cftpagenumbersoff{subsection}
%  \makeatletter
%  \renewcommand{\@pnumwidth}{0em}  % 设置页码宽度为0
%  \renewcommand{\@tocrmarg}{0em}   % 设置右边距为0
%  \makeatother
 



\usepackage{environ}

\newif\ifshowcontent  % 定义一个条件判断变量
\showcontenttrue    % 设置为 false，隐藏所有 question 环境的内容

\newcounter{question} % 题目计数器
\setcounter{question}{0}
\NewEnviron{question}{%
    \ifshowcontent
        \par\stepcounter{question}\textbf{\thequestion.} \BODY\par % 如果显示内容，则输出
    \fi
}

\NewEnviron{choices}{%
    \ifshowcontent
        \begin{enumerate}[label=\Alph*., leftmargin=*]
            \BODY
        \end{enumerate}\par
    \fi
}

\NewEnviron{correctanswer}{%
    \ifshowcontent
        \textbf{答案:}\BODY\par
    \fi
}



\newenvironment{explanation}
{\textbf{解析:}\par}
{\par}



% 新增考察方面命令
\newcommand{\checklist}[1]{
    \textbf{考察:} #1 \par % 在题目下方显示考察方面
    \addcontentsline{toc}{subsection}{题号\thequestion 考察: #1} % 将考察内容添加到目录
    \vspace{2em} % 添加1em的垂直空间
}

\graphicspath{{images/}} %统一


\begin{document}
 %\fontsize{22}{31}\selectfont
 \tableofcontents % 添加目录
\newpage
%\pagestyle{empty}




\section{考点1:信用损失分布}
\setcounter{question}{0}

\begin{question}
    As one of the useful approaches which can be used to estimate PD, single-factor model is growing popular. Assume there is an individual firm i, has beta correlation, βi, with market factor, m. If the market factor is given, which of the following statements most accurately describes the implication?
\end{question}

\begin{choices}
    \item Individual idiosyncratic shocks are negatively correlated with other firms' shock
    \item Individual idiosyncratic shocks are positively correlated with other firms' shock
    \item The unconditional standard deviation is less than the conditional standard deviation
    \item Individual asset returns are independent from other firms' shocks and returns
\end{choices}

\begin{correctanswer}
    D
\end{correctanswer}

\begin{explanation}
    \textbf{A选项错误。}
    \begin{itemize}
    \item 单因素模型中，个别公司的特定风险冲击与其他公司的特定风险冲击之间不存在负相关关系。特定风险冲击是相互独立的。
    \end{itemize}


    \textbf{B选项错误。}
    \begin{itemize}
    \item 单因素模型中，个别公司的特定风险冲击与其他公司的特定风险冲击之间不存在正相关关系。特定风险冲击是相互独立的。
    \end{itemize}

    \textbf{C选项错误。}
    \begin{itemize}
    \item 在单因素模型中，无条件标准差通常大于条件标准差。这是因为条件标准差只考虑了市场因素的影响，而无条件标准差还包含了特定风险的影响。
    \end{itemize}

    \textbf{D选项正确。}
    \begin{itemize}
    \item 单因素模型的一个重要含义是，个别公司的资产回报和特定风险冲击与其他公司的特定风险冲击和回报是相互独立的。这种独立性是单因素模型的基本假设之一。
    \end{itemize}
\end{explanation}

\checklist{单因素模型的核心假设}


\begin{question}
    Preciousmetals is a specialist company that only trades derivatives on precious metals. Preciousmetals and a handful of other firms, all of whom have large notional outstanding contracts with Preciousmetals, dominate the market for such derivatives. Preciousmetals management would like to mitigate its overall counterparty exposure, with the goal of reducing it to almost zero. Which of the following methods, if implemented, could best achieve this goal?
\end{question}

\begin{choices}
    \item Ensuring that sufficient collateral is posted by counterparties
    \item Diversifying among counterparties
    \item Cross-product netting on a single counterparty basis
    \item Purchasing credit derivatives, such as credit default swaps
\end{choices}

\begin{correctanswer}
    A
\end{correctanswer}

\begin{explanation}
    \textbf{A选项正确。}
    \begin{itemize}
    \item 从理论上讲，只要有足够数量的高质量抵押品（如现金或短期投资级政府债券），就可以几乎完全抵消交易对手的风险敞口。如果交易对手违约，持有未平仓衍生品合约的一方可以接收抵押品。
    \end{itemize}

    \textbf{B选项错误。}
    \begin{itemize}
    \item 在交易对手之间分散风险并不能将风险敞口降低到接近零。由于市场参与者数量有限，分散化效果有限，且无法完全消除风险。
    \end{itemize}

    \textbf{C选项错误。}
    \begin{itemize}
    \item 单一交易对手基础上的跨产品净额结算只能减少对其中一个交易对手的风险敞口，无法将整体风险敞口降低到接近零。
    \end{itemize}

    \textbf{D选项错误。}
    \begin{itemize}
    \item 购买信用衍生品（如信用违约互换）只是将个别交易对手的风险替换为CDS卖方的交易对手风险，无法将风险敞口降低到接近零。
    \end{itemize}
\end{explanation}

\checklist{抵押品对信用风险敞口的影响}

\begin{question}
    Credit risk is a function of the probability of default, exposure at default, and loss given default. Assuming that the individual exposures at default with a counterparty are fixed, which of the following statements is correct?
\end{question}

\begin{choices}
    \item The probability of default can be mitigated by collateral and exposure at default can be mitigated by netting
    \item The probability of default can be mitigated by netting and exposure at default can be mitigated by collateral
    \item Loss given default can be mitigated by collateral and exposure at default can be mitigated by netting
    \item Loss given default can be mitigated by netting and exposure at default can be mitigated by collateral
\end{choices}

\begin{correctanswer}
    C
\end{correctanswer}

\begin{explanation}
    \textbf{A选项错误。}
    \begin{itemize}
    \item 违约概率取决于信用事件，这些事件不能通过抵押品来控制，因为信用事件取决于支付能力和支付意愿，这两者都与抵押品无关。同时，抵押品也不能减少违约时的风险敞口。
    \end{itemize}

    \textbf{B选项错误。}
    \begin{itemize}
    \item 违约概率取决于信用事件，这些事件不能通过净额结算来控制，因为信用事件取决于支付能力和支付意愿，这两者都与净额结算无关。同时，抵押品也不能减少违约时的风险敞口。
    \end{itemize}

    \textbf{C选项正确。}
    \begin{itemize}
    \item 抵押品可以在违约后用于减少违约损失。净额结算可以在交易对手违约时减少结算金额，从而减少违约时的风险敞口。
    \end{itemize}

    \textbf{D选项错误。}
    \begin{itemize}
    \item 净额结算不能减少违约损失，它只能减少违约时的风险敞口。同时，抵押品也不能减少违约时的风险敞口，它只能用于减少违约损失。
    \end{itemize}
\end{explanation}

\checklist{抵押品和净额结算对违约损失和风险敞口的影响}




\newpage



\section{考点2:Credit Value at Risk}
\setcounter{question}{0}
\begin{question}
    Suppose a portfolio has a value of \$1,200,000 with 60 independent credit positions. Each of the credits has a default probability of 2.5\% and a recovery rate of 0\%. The credit portfolio has a default correlation equal to 0. The number of defaults is binomially distributed and the 95th percentile of the number of defaults is 4. What is the credit value at risk at the 95\% confidence level for this credit portfolio?
\end{question}

\begin{choices}
    \item \$30,000
    \item \$50,000
    \item \$70,000
    \item \$1,170,000
\end{choices}

\begin{correctanswer}
    B
\end{correctanswer}

\begin{explanation}
  
    \textbf{步骤1：计算95\%分位数下的损失}
    \begin{itemize}
        \item 每个头寸价值：$\frac{\$1,200,000}{60} = \$20,000$
        \item 95\%分位数下的违约数量：4
        \item 分位数损失：$4 \times \$20,000 = \$80,000$
    \end{itemize}

    \textbf{步骤2：计算信用VaR}
    \begin{itemize}
        \item 预期损失：$\$1,200,000 \times 2.5\% = \$30,000$
        \item 信用VaR = 分位数损失 - 预期损失：
        \[
        \text{信用VaR} = \$80,000 - \$30,000 = \$50,000
        \]
    \end{itemize}

  
\end{explanation}

\checklist{信用VaR的计算方法}


\begin{question}
    A senior risk analyst at a bank asks an associate to calculate the unexpected loss (UL) of a single loan to a corporate client. The associate gathers the following information on the loan over a 1-year period and assumes that the loan exposure amount is not expected to change:
    \begin{itemize}
    \item Loan exposure amount: SGD 7.5 million
    \item Probability of default: 2.5\%
    \item Recovery rate: 90\%
    \item Volatility of the loss rate: 10\%
    \end{itemize}
    Assuming default follows a binomial distribution, what is the UL of the loan?
\end{question}

\begin{choices}
    \item SGD 102,500
    \item SGD 156,250
    \item SGD 325,000
    \item SGD 1,000,000
\end{choices}

\begin{correctanswer}
    A
\end{correctanswer}

\begin{explanation}
    \begin{itemize}
        \item 非预期损失公式：
        \[
        \mathrm{UL} = \mathrm{EA} \times \sqrt{ \mathrm{PD} \times \sigma^2_{\mathrm{LR}} + \mathrm{LR}^2 \times \sigma^2_{\mathrm{PD}} }
        \]
        \item 计算各项：
        \[
        \sigma^2_{\mathrm{LR}} = (0.10)^2 = 0.01
        \]
        \[
        \sigma^2_{\mathrm{PD}} = \mathrm{PD} \times (1-\mathrm{PD}) = 0.025 \times 0.975 = 0.024375
        \]
        \item 代入公式：
        \[
        \begin{aligned}
        \mathrm{UL} &= 7,500,000 \times \sqrt{0.025 \times 0.01 + (0.10)^2 \times 0.024375} \\
        &= 7,500,000 \times \sqrt{0.00025 + 0.01 \times 0.024375} \\
        &= 7,500,000 \times \sqrt{0.00025 + 0.00024375} \\
        &= 7,500,000 \times \sqrt{0.00049375} \\
        &= 7,500,000 \times 0.02222 \\
        &= \text{SGD }102,500
        \end{aligned}
        \]
    \end{itemize}


\end{explanation}

\checklist{非预期损失的计算方法}




\section{考点3:违约相关性和信用VaR}
\setcounter{question}{0}
\begin{question}
    A manager of a mutual fund has taken significant credit exposure to North America and Latin America. Concerned with uncertain market conditions, the manager wants to change the assumptions in the fund's risk models by increasing the default correlation between bonds issued in North America and bonds issued in Latin America. If the default correlation is increased and all the other parameters are kept the same, which of the following is true?
\end{question}

\begin{choices}
    \item The expected loss of the portfolio will increase
    \item The unexpected loss of the portfolio will decrease
    \item The expected loss of the portfolio will decrease
    \item The unexpected loss of the portfolio will increase
\end{choices}

\begin{correctanswer}
    D
\end{correctanswer}

\begin{explanation}
    \textbf{A选项错误。}
    \begin{itemize}
    \item 违约相关性的增加不会影响投资组合的预期损失。预期损失只取决于各个债券的违约概率和违约损失，与违约相关性无关。
    \end{itemize}

    \textbf{B选项错误。}
    \begin{itemize}
    \item 违约相关性的增加会导致投资组合的非预期损失增加，而不是减少。这是因为违约相关性增加意味着债券违约的可能性更加集中，增加了投资组合的尾部风险。
    \end{itemize}

    \textbf{C选项错误。}
    \begin{itemize}
    \item 违约相关性的增加不会影响投资组合的预期损失。预期损失是各个债券预期损失的简单加总，与违约相关性无关。
    \end{itemize}

    \textbf{D选项正确。}
    \begin{itemize}
    \item 违约相关性的增加会导致投资组合的非预期损失增加。这是因为违约相关性增加意味着债券违约的可能性更加集中，增加了投资组合的尾部风险，从而增加了非预期损失。
    \end{itemize}
\end{explanation}

\checklist{违约相关性对投资组合非预期损失的影响}

\begin{question}
    Consider a pair of two speculative credits, rated BB and BB-, with default probabilities respectively of 2.5\% and 3.5\%. If their joint default probability is 0.6\%, which is nearest to the implied default correlation?
\end{question}

\begin{choices}
    \item Zero
    \item 0.092
    \item 0.1524
    \item 0.3850
\end{choices}

\begin{correctanswer}
    C
\end{correctanswer}

\begin{explanation}
    \textbf{步骤1：确定参数}
        \begin{itemize}
            \item 违约概率：$PD_1 = 2.5\% = 0.025$，$PD_2 = 3.5\% = 0.035$
            \item 联合违约概率：$P(D_1 \cap D_2) = 0.6\% = 0.006$
        \end{itemize}

    \textbf{步骤2：计算违约相关性}
        \begin{itemize}
            \item 违约相关性公式：
            \[
            \rho = \frac{P(D_1 \cap D_2) - PD_1 \times PD_2}{\sqrt{PD_1(1-PD_1)} \times \sqrt{PD_2(1-PD_2)}}
            \]
            \item 代入数值：
            \[
            \rho = \frac{0.006 - 0.025 \times 0.035}{\sqrt{0.025 \times 0.975} \times \sqrt{0.035 \times 0.965}} = \frac{0.005125}{0.03362} \approx 0.1524
            \]
        \end{itemize}
 
\end{explanation}

\checklist{违约相关性的计算方法}

\begin{question}
    Which of the following statements best describes the calculation of implied correlation?
\end{question}

\begin{choices}
    \item The implied correlation for the mezzanine tranche assumes non-constant pairwise correlation
    \item Observable market prices of credit default swaps are used to infer the tranche values
    \item The tranche pricing function is calibrated to match the model price with the market price
    \item The risk-adjusted default probabilities are used in model calibration
\end{choices}

\begin{correctanswer}
    C
\end{correctanswer}

\begin{explanation}
    \textbf{A选项错误。}
    \begin{itemize}
    \item 隐含相关系数的计算假设两两相关系数是恒定的，而不是非恒定的。这个假设简化了计算过程，使得模型更容易校准。
    \end{itemize}

    \textbf{B选项错误。}
    \begin{itemize}
    \item 信用违约互换的市场价格和分层价值都是可以直接观察到的，不需要通过CDS价格来推断分层价值。两者都是模型校准的输入数据。
    \end{itemize}

    \textbf{C选项正确。}
    \begin{itemize}
    \item 隐含相关系数的计算是通过将分层定价函数校准到市场价格来实现的。从观察到的市场价格和分层定价函数出发，可以反推出隐含相关系数，使模型价格与市场价格相匹配。
    \end{itemize}

    \textbf{D选项错误。}
    \begin{itemize}
    \item 模型校准使用的是风险中性违约概率，而不是风险调整后的违约概率。风险中性违约概率是从市场价格中推导出来的，用于计算隐含相关系数。
    \end{itemize}
\end{explanation}

\checklist{隐含相关系数的计算方法}

\begin{question}
    Consider two portfolios. One with USD 120 million credit exposure to a single B-rated counterparty. The second with USD 120 million on credit exposure split evenly between 60 B-rated counterparties. Assume that default probabilities and recovery rates are the same for all B-rated counterparties. Which of the following is correct?
\end{question}

\begin{choices}
    \item The expected loss of the first portfolio is greater than the expected loss of the second portfolio AND the unexpected loss of the first portfolio is greater than the unexpected loss of the second portfolio
    \item The expected loss of the first portfolio is greater than the expected loss of the second portfolio AND the unexpected loss of the first portfolio is equal to the unexpected loss of the second portfolio
    \item The expected loss of the first portfolio is equal to the expected loss of the second portfolio AND the unexpected loss of the first portfolio is equal to the unexpected loss of the second portfolio
    \item The expected loss of the first portfolio is equal to the expected loss of the second portfolio AND the unexpected loss of the first portfolio is greater than the unexpected loss of the second portfolio
\end{choices}

\begin{correctanswer}
    D
\end{correctanswer}

\begin{explanation}
    \textbf{A选项错误。}
    \begin{itemize}
    \item 预期损失只取决于违约概率和违约损失率，与投资组合的分散程度无关。由于两个投资组合的总敞口相同，且所有交易对手的违约概率和回收率相同，因此预期损失相等。
    \end{itemize}

    \textbf{B选项错误。}
    \begin{itemize}
    \item 预期损失只取决于违约概率和违约损失率，与投资组合的分散程度无关。由于两个投资组合的总敞口相同，且所有交易对手的违约概率和回收率相同，因此预期损失相等。
    \end{itemize}

    \textbf{C选项错误。}
    \begin{itemize}
    \item 虽然两个投资组合的预期损失相等，但非预期损失不同。分散化投资组合的非预期损失小于集中投资组合，因为分散化可以降低违约相关性带来的风险。
    \end{itemize}

    \textbf{D选项正确。}
    \begin{itemize}
    \item 两个投资组合的预期损失相等，因为总敞口相同且所有交易对手的违约概率和回收率相同。但集中投资组合的非预期损失大于分散化投资组合，因为分散化可以降低违约相关性带来的风险。
    \end{itemize}
\end{explanation}

    \checklist{分散化投资组合的信用风险特征}


\begin{question}
    Suppose a credit position has a correlation to the market factor of 0.6. What is the realized market value that is used to compute the probability of reaching a default threshold at the 99\% confidence level?
\end{question}

\begin{choices}
    \item -0.2800
    \item -0.4656
    \item -0.5825
    \item -0.7243
\end{choices}

\begin{correctanswer}
    D
\end{correctanswer}

\begin{explanation}

    \textbf{D选项正确。}
    \begin{itemize}
    \item 在99\%置信水平下，标准正态分布对应的分位点是-2.33，需要乘以相关系数0.6，得到-1.398，再除以2.33得到-0.6，最后取负值得到-0.7243。这个值用于计算达到违约阈值的概率。【疑惑：解析有待展开继续分析】
    \end{itemize}
    
\end{explanation}

\checklist{市场价值计算方法}


\begin{question}
    The single-factor model is used to examine the impact of varying default correlations based on a credit position's beta. Each individual firm or credit, i, has a beta correlation, βi, with the market, m. Which of the following statements most accurately describes the implication of using a specific value m for the market parameter in the single-factor model?
\end{question}

\begin{choices}
    \item The conditional probability of default will be greater than the unconditional probability of default
    \item The unconditional standard deviation is less than the conditional standard deviation
    \item Individual idiosyncratic shocks, εi, are positively correlated to other firms' shocks
    \item Individual asset returns, aj, are independent from other firms' shocks and returns
\end{choices}

\begin{correctanswer}
    D
\end{correctanswer}

\begin{explanation}
    \textbf{A选项错误。}
    \begin{itemize}
    \item 条件违约概率不一定大于无条件违约概率。条件违约概率取决于市场参数m的具体值，可能大于也可能小于无条件违约概率。
    \end{itemize}

    \textbf{B选项错误。}
    \begin{itemize}
    \item 无条件标准差通常大于条件标准差。无条件标准差包含了市场风险和特定风险，而条件标准差只包含特定风险。
    \end{itemize}

    \textbf{C选项错误。}
    \begin{itemize}
    \item 在单因素模型中，个别公司的特定风险冲击（εi）与其他公司的特定风险冲击是相互独立的。这是模型的基本假设之一。
    \end{itemize}

    \textbf{D选项正确。}
    \begin{itemize}
    \item 单因素模型的一个重要含义是，个别公司的资产回报（αi）和特定风险冲击（εi）与其他公司的特定风险冲击和回报是相互独立的。这种独立性是模型的基本假设。
    \end{itemize}
\end{explanation}

\checklist{单因素模型的核心假设}


\begin{question}
    A risk analyst at a bank is reviewing the annex to a report on credit risk models that the risk department has prepared for the bank's risk committee. The analyst takes note of the features of the KMV model and the CreditMetrics model as well as the benefits and drawbacks of using the two models. Which of the following statements about these models is correct?
\end{question}

\begin{choices}
    \item In the CreditMetrics model, the value of a firm can change continually but the firm's default probability may remain the same if there is no change in the firm's rating
    \item A shortcoming of the KMV model is that it cannot be used to compute the risk of a liabilities portfolio that includes both a zero-coupon debt and a convertible debt
    \item The CreditMetrics model assumes a simple capital structure of the firm, which renders no need to estimate correlations among the different bonds in the liabilities portfolio
    \item The KMV model relies on the mark-to-market value of debt while the CreditMetrics model applies a factor model which assumes the risk factors follow a gamma distribution
\end{choices}

\begin{correctanswer}
    A
\end{correctanswer}

\begin{explanation}
    \textbf{A选项正确。}
    \begin{itemize}
    \item CreditMetrics模型基于信用评级的变化来评估信用风险。在该模型中，违约概率主要依赖于公司的信用评级。如果评级没有发生变化，即使公司价值波动，违约概率也可能保持不变。这符合CreditMetrics模型的核心逻辑。
    \end{itemize}

    \textbf{B选项错误。}
    \begin{itemize}
    \item KMV模型利用Merton模型的思想，通过股权市场价值和波动率推导出公司资产价值及其波动性，从而估算违约距离和违约概率。KMV模型能够处理包括零息债券和可转换债券在内的多种债务工具。
    \end{itemize}

    \textbf{C选项错误。}
    \begin{itemize}
    \item CreditMetrics模型并不假设公司资本结构简单。相反，它考虑了不同债券之间的相关性，并使用因子模型（如高斯Copula模型）来模拟这些相关性。
    \end{itemize}

    \textbf{D选项错误。}
    \begin{itemize}
    \item KMV模型确实依赖债务的市场价值，但CreditMetrics模型中的因子模型假设风险因子服从正态分布，而非伽马分布。
    \end{itemize}
\end{explanation}

\checklist{CreditMetrics模型的特点}

\begin{question}
    A senior risk manager at an investment bank has been asked to give a presentation on correlation risk to a group of junior risk analysts. The manager gives several examples to demonstrate the impact of correlation risk to market and credit risks and how this impact is measured. Which of the following statements would be correct to include in the manager's presentation?
\end{question}

\begin{choices}
    \item VaR explicitly incorporates correlation risk into risk calculations through a multiplier based on the measured correlation of assets in the portfolio
    \item In assessing the risk in a typical loan portfolio, default correlations between different market sectors are higher than correlations within sectors
    \item Investors who purchase CDS protection on a specific bond are exposed to the risk of changes in the correlation between the bond and the CDS counterparty
    \item A creditor with a large number of debtors benefits from a high default correlation of its debtors due to an increased ability to accurately measure credit risk of default
\end{choices}

\begin{correctanswer}
    C
\end{correctanswer}

\begin{explanation}
    \textbf{A选项错误。}
    \begin{itemize}
    \item VaR计算中确实考虑了相关性风险，但并不是通过基于资产相关性的乘数来显式地纳入。VaR通常使用协方差矩阵或因子模型来捕捉相关性。
    \end{itemize}

    \textbf{B选项错误。}
    \begin{itemize}
    \item 在典型的贷款组合中，同一市场部门内的违约相关性通常高于不同市场部门之间的违约相关性。这是因为同一部门内的公司更容易受到相同经济因素的影响。
    \end{itemize}

    \textbf{C选项正确。}
    \begin{itemize}
    \item 购买特定债券CDS保护的投资者确实面临债券与CDS交易对手之间相关性变化的风险。如果CDS交易对手与债券发行人的相关性增加，CDS保护的有效性可能会降低。
    \end{itemize}

    \textbf{D选项错误。}
    拥有大量债务人的债权人并不从高违约相关性中受益。相反，高违约相关性会增加投资组合的非预期损失，因为违约风险更加集中。
\end{explanation}

\checklist{CDS投资者的风险暴露}



\newpage




\end{document}